\begin{abstract}
Vision-Language-Action (VLA) models increasingly rely on action chunking---predicting a sequence of future actions in a single forward pass---to improve execution efficiency and temporal coherence. However, action chunking introduces a systematic discontinuity at chunk boundaries: the last action of one chunk and the first action of the next are predicted independently, often resulting in abrupt velocity jumps that degrade manipulation performance. In this paper, we present TAR (Temporal Action Regularization), a lightweight training-time regularization that penalizes high velocities at the start and end of each predicted action chunk. TAR requires only a single additional loss term and incurs zero inference overhead. We integrate TAR into BitVLA, a 1-bit VLA model (3\,B parameters, 1.4\,GB memory), and evaluate on the LIBERO-10 benchmark. TAR improves the task success rate from 86.4\% (baseline) to 93.6\% (+7.2\,pp) with $\lambda_{\text{TAR}}=0.1$, while simultaneously reducing intra-chunk velocity variation (S1) by 36\%, peak action jumps (MaxJump) by 53\%, and gripper switching noise by 24\%. Comprehensive ablation over four $\lambda$ values (0, 0.01, 0.05, 0.10) reveals a monotonically increasing trend in both success rate and action smoothness, confirming that TAR functions as both a smoothing regularizer and a generalization-enhancing constraint with no observed performance saturation in the tested range.
\end{abstract}

\begin{IEEEkeywords}
Vision-Language-Action Models, Action Chunking, Temporal Regularization, Robot Manipulation, LIBERO
\end{IEEEkeywords}
